\documentclass[UTF8,zihao=-4,oneside,openany]{ctexbook}
\usepackage{amsmath,amssymb,amsthm,bm,geometry,booktabs,enumitem,xcolor}
\usepackage[unicode,hidelinks]{hyperref}
\geometry{a4paper,margin=2.4cm}
\setlist{itemsep=2pt,topsep=4pt}
\setcounter{tocdepth}{2}
\setcounter{secnumdepth}{2}
\allowdisplaybreaks[2]
\newcommand{\step}[1]{\par\medskip\noindent\textbf{#1}\par\nopagebreak}
\newcommand{\ind}{\mathbf{1}}
\newcommand{\Var}{\operatorname{Var}}
\newcommand{\MSE}{\operatorname{MSE}}
\newcommand{\Be}{\operatorname{Be}}
\newcommand{\Ga}{\operatorname{Ga}}
\newcommand{\IGa}{\operatorname{IGa}}
\newcommand{\Bin}{\operatorname{Bin}}
\newcommand{\BF}{\operatorname{BF}}
\emergencystretch=2em
\title{贝叶斯统计课堂笔记}
\author{根据课堂板书与课程讲义整理}
\date{2026年10月27日修订\\绪论、第一章与第二章（至2.6似然原理）}
\begin{document}
\hypersetup{pageanchor=false}
\maketitle
\clearpage
\pagenumbering{roman}
\hypersetup{pageanchor=true}
\tableofcontents
\chapter*{整理说明}
\addcontentsline{toc}{chapter}{整理说明}
本笔记按章节整理，公式、记号和主要推导以所提供的课堂板书为主，课件用于核对题目及补足推导。第二章例题编号统一采用课件编号；例如板书的“例2.4”“例2.5”“例2.10”“例2.11”，分别对应课件的例2.2.3、2.2.4、2.4.1、2.4.2。

计算题尽量依照“似然函数、先验分布、后验核、归一化或识别分布、所求结论”的顺序展开。区间估计、检验和预测在得到后验之后再继续计算。例2.2.6的完整推导依据课件补入；正态预测中的积分配方也参考课件展开。板书与课件中的个别近似或数值不一致处，在相应例题中注明。

本次60张照片与笔记内容的逐项对照另存于同目录的《板书对照.md》，其中照片编号按本次上传的先后顺序排列。旧版文件已另作备份。
\chapter*{绪论}
\addcontentsline{toc}{chapter}{绪论}
\section*{经典统计与贝叶斯统计}
统计推断通常利用总体信息（总体分布或分布族）、样本信息（观测数据）和先验信息（观测前关于未知参数的知识）。经典统计主要依据总体与样本信息；贝叶斯统计将未知参数视为随机变量，以先验分布表示先验信息，再用样本信息更新为后验分布。

概率解释的差异在于：频率学派将概率理解为重复试验下的长期频率，参数通常视为固定常数；贝叶斯学派允许概率表示对未知事件或参数的信念，参数以概率分布刻画。贝叶斯推断的基本流程为
\[
\text{先验信息}+\text{样本信息}\xrightarrow{\text{贝叶斯公式}}\text{后验信息}\xrightarrow{}\text{推断或决策}.
\]
经典方法的区间覆盖率是重复抽样下的性质；贝叶斯可信区间则可直接表述为给定数据后参数落入区间的后验概率。

\section*{医疗试验引例}
课堂引例设控制饮食的治愈率为 $0.1$。隐藏因素分为 $F_A,F_B$，且
\[
P(F_A)=0.9,\quad P(F_B)=0.1,\quad P(\text{药物有效}\mid F_A)=0,
\quad P(\text{药物有效}\mid F_B)=0.95.
\]
因而治疗组个体治愈率为
\[
P(\text{治愈}\mid\text{drug})=0.9(0.1)+0.1\{1-(1-0.1)(1-0.95)\}=0.1855.
\]
即模型给出的治疗组治愈率为 $18.55\%$，相对 $10\%$ 的饮食治愈率提高 $8.55$ 个百分点。这个例子用于说明群体频率比较与针对个体的概率判断所回答的问题不同。

\section*{事件形式的贝叶斯公式}
若 $A_1,\ldots,A_k$ 构成完备事件组，且 $P(B)>0$，则
\[
P(A_j\mid B)=\frac{P(B\mid A_j)P(A_j)}{\sum_{i=1}^kP(B\mid A_i)P(A_i)}.
\]
分子是“假设 $A_j$ 成立且观察到 $B$”的概率，分母是观察到 $B$ 的总概率。

\subsection*{医学检测}
令 $A$ 表示患病、$B$ 表示检测阳性。即使检测灵敏度 $P(B\mid A)$ 很高，阳性后的患病概率仍为
\[
P(A\mid B)=\frac{P(B\mid A)P(A)}{P(B\mid A)P(A)+P(B\mid A^c)P(A^c)}.
\]
必须同时考虑患病率和假阳性率，不能把 $P(B\mid A)$ 与 $P(A\mid B)$ 混淆。

\section*{参数的密度形式与共轭更新}
若样本 $x$ 在给定参数 $\theta$ 下的似然为 $L(x\mid\theta)$，先验密度为 $\pi(\theta)$，则
\[
\pi(\theta\mid x)=\frac{L(x\mid\theta)\pi(\theta)}{m(x)},\qquad
m(x)=\int_\Theta L(x\mid\theta)\pi(\theta)\,d\theta.
\]
其中 $m(x)$ 是边缘似然（证据），保证后验密度积分为 $1$。计算时常先写
\[
\pi(\theta\mid x)\propto L(x\mid\theta)\pi(\theta),
\]
然后辨认分布核并补上归一化常数。

\subsection*{二项--贝塔示例}
若 $X\mid\theta\sim\operatorname{Bin}(n,\theta)$，观测成功数为 $x$，则
\[
L(x\mid\theta)=\binom{n}{x}\theta^x(1-\theta)^{n-x}\propto\theta^x(1-\theta)^{n-x}.
\]
取均匀先验 $\theta\sim\operatorname{Be}(1,1)$，有
\[
\pi(\theta\mid x)=(n+1)\binom{n}{x}\theta^x(1-\theta)^{n-x},\quad
\theta\mid x\sim\operatorname{Be}(x+1,n-x+1).
\]
边缘概率为 $m(x)=1/(n+1)$（$x=0,1,\ldots,n$）。一般地，若先验为 $\operatorname{Be}(\alpha,\beta)$，则后验为 $\operatorname{Be}(\alpha+x,\beta+n-x)$。

\section*{贝叶斯统计的发展与应用}
Thomas Bayes 的工作于 1763 年由 Richard Price 在英国皇家学会宣读发表；Laplace 对贝叶斯推理的发展作出重要贡献。统计决策理论的发展、决策问题对先验知识的需要，以及经济、工业、医学等应用推动了贝叶斯统计的发展。

贝叶斯方法的重要挑战包括先验分布的确定和后验分布的计算。课程介绍了无信息先验、共轭先验、最大熵以及自助法和随机加权法等先验确定思路。现代数值近似和 MCMC 方法提高了复杂后验的可计算性。应用领域包括经济金融与保险、生物医学与生态、可靠性和机器学习等。课程还提及客观贝叶斯、主观贝叶斯、稳健贝叶斯、频率贝叶斯与拟贝叶斯等方向。

\clearpage
\pagenumbering{arabic}
\chapter{先验分布与后验分布}
\section{三种信息与先验分布}
总体信息是总体分布或所属分布族提供的信息；样本信息来自抽取的数据；先验信息是观测前关于未知量的经验、历史资料或已有认识。贝叶斯统计综合三者。

把未知参数 $\theta$ 看作随机变量，用概率分布描述观测前对其不确定性的认识，称为先验分布，记作 $\pi(\theta)$。先验概率可表达主观信念，不必是重复试验频率。离散参数用先验概率质量函数，连续参数用先验密度。

\section{贝叶斯公式}
\subsection{事件形式}
若 $A_1,\dots,A_k$ 为完备事件组，$P(B)>0$，则
\[
P(A_j\mid B)=\frac{P(B\mid A_j)P(A_j)}{\sum_iP(B\mid A_i)P(A_i)}.
\]
\subsection{密度形式}
对连续参数，样本 $x$ 的联合密度记 $h(x,\theta)=L(x\mid\theta)\pi(\theta)$，样本边缘密度
\[
m(x)=\int_\Theta h(x,\theta)\,d\theta,
\]
故
\[
\pi(\theta\mid x)=\frac{h(x,\theta)}{m(x)}
=\frac{L(x\mid\theta)\pi(\theta)}{\int_\Theta L(x\mid t)\pi(t)\,dt}.
\]
注意 $L(x\mid\theta)$ 是观测数据固定后关于参数的函数；$\pi(\theta)$ 是先验；分母不依赖 $\theta$，用于归一化。

\subsection{二项模型与均匀先验}
若 $X\mid\theta\sim\operatorname{Bin}(n,\theta)$，则
\[
L(x\mid\theta)=\binom{n}{x}\theta^x(1-\theta)^{n-x},\quad x=0,1,\ldots,n.
\]
取 $\theta\sim U(0,1)=\operatorname{Be}(1,1)$，则
\[
m(x)=\int_0^1\binom{n}{x}\theta^x(1-\theta)^{n-x}\,d\theta=\frac1{n+1},
\]
以及
\[
\theta\mid x\sim\operatorname{Be}(x+1,n-x+1).
\]
其密度为 $\frac{\Gamma(n+2)}{\Gamma(x+1)\Gamma(n-x+1)}\theta^x(1-\theta)^{n-x}$。

\section{共轭先验分布}
若先验与似然相乘后所得后验仍属于同一分布族，则该先验称为该模型参数的共轭先验。共轭性是针对“指定模型中的指定参数”而言。优点是后验计算方便、超参数常可解释；选择先验时合理性优先于计算便利。

\subsection{正态均值，方差已知}
设 $X_1,\ldots,X_n\mid\theta\overset{iid}{\sim}N(\theta,\sigma^2)$，$\sigma^2$ 已知，先验 $\theta\sim N(\mu,\tau^2)$。令 $\bar x=n^{-1}\sum_i x_i$，则
\[
L(x\mid\theta)\propto\exp\left[-\frac{n}{2\sigma^2}(\theta-\bar x)^2\right],\qquad
\pi(\theta)\propto\exp\left[-\frac1{2\tau^2}(\theta-\mu)^2\right].
\]
相乘并配方得 $\theta\mid x\sim N(\mu_1,\tau_1^2)$，其中
\[
\frac1{\tau_1^2}=\frac n{\sigma^2}+\frac1{\tau^2},\qquad
\mu_1=\frac{n\bar x/\sigma^2+\mu/\tau^2}{n/\sigma^2+1/\tau^2}
=\gamma\bar x+(1-\gamma)\mu,
\quad \gamma=\frac{n/\sigma^2}{n/\sigma^2+1/\tau^2}.
\]
后验均值是样本均值与先验均值的精度加权平均；后验精度等于样本均值分布精度与先验精度之和。

\subsection{二项概率，贝塔先验}
若 $X\mid\theta\sim\operatorname{Bin}(n,\theta)$，且 $\theta\sim\operatorname{Be}(\alpha,\beta)$，则
\[
\pi(\theta\mid x)\propto\theta^{\alpha+x-1}(1-\theta)^{\beta+n-x-1},\quad
\theta\mid x\sim\operatorname{Be}(\alpha+x,\beta+n-x).
\]
因此
\[
E(\theta\mid x)=\frac{\alpha+x}{\alpha+\beta+n}
=\gamma\frac xn+(1-\gamma)\frac\alpha{\alpha+\beta},\quad
\gamma=\frac n{\alpha+\beta+n},
\]
\[
\operatorname{Var}(\theta\mid x)=
\frac{(\alpha+x)(\beta+n-x)}{(\alpha+\beta+n)^2(\alpha+\beta+n+1)}
=\frac{E(\theta\mid x)[1-E(\theta\mid x)]}{\alpha+\beta+n+1}.
\]
固定超参数且 $n$ 较大时，后验均值主要由 $x/n$ 决定，后验方差趋于零；这不表示任意增加一个新观测，后验方差都必然单调下降。

\subsection{Gamma、Beta 函数及常用分布}
Gamma 函数与 Beta 函数定义为
\[
\Gamma(a)=\int_0^\infty x^{a-1}e^{-x}\,dx,\quad a>0,
\qquad
B(a,b)=\int_0^1t^{a-1}(1-t)^{b-1}\,dt
=\frac{\Gamma(a)\Gamma(b)}{\Gamma(a+b)}.
\]
常用关系为 $\Gamma(a+1)=a\Gamma(a)$、$\Gamma(n+1)=n!$。若 $X\sim\operatorname{Be}(a,b)$，则
\[
f(x)=\frac{x^{a-1}(1-x)^{b-1}}{B(a,b)},\quad 0<x<1,
\quad E(X)=\frac a{a+b},\quad
\operatorname{Var}(X)=\frac{ab}{(a+b)^2(a+b+1)}.
\]
若 $X_1\sim\operatorname{Ga}(a_1,\lambda)$、$X_2\sim\operatorname{Ga}(a_2,\lambda)$ 相互独立且率参数相同，则 $X_1+X_2\sim\operatorname{Ga}(a_1+a_2,\lambda)$。指数分布是 $a=1$ 的 Gamma 分布。

本笔记按课堂板书采用 Gamma 的率参数化：若 $X\sim\operatorname{Ga}(\alpha,\lambda)$，则
\[
f(x)=\frac{\lambda^\alpha}{\Gamma(\alpha)}x^{\alpha-1}e^{-\lambda x},\quad x>0;
\quad E(X)=\frac\alpha\lambda,\quad \operatorname{Var}(X)=\frac\alpha{\lambda^2}.
\]
若 $Y=X^{-1}$，变量变换给出
\[
f_Y(y)=\frac{\lambda^\alpha}{\Gamma(\alpha)}y^{-\alpha-1}e^{-\lambda/y},\quad y>0,
\]
记 $Y\sim\operatorname{IGa}(\alpha,\lambda)$。当 $\alpha>1$ 时 $E(Y)=\lambda/(\alpha-1)$。

\subsection{例1.3.3：正态方差的倒 Gamma 共轭先验}
设 $X_i\mid\sigma^2\overset{iid}{\sim}N(\mu,\sigma^2)$，$\mu$ 已知，令 $T=\sum_{i=1}^n(x_i-\mu)^2$。
\step{Step 1：似然函数}
\[
L(\sigma^2;x)=\prod_{i=1}^n\frac1{\sqrt{2\pi\sigma^2}}
\exp\left[-\frac{(x_i-\mu)^2}{2\sigma^2}\right]
=(2\pi\sigma^2)^{-n/2}e^{-T/(2\sigma^2)}.
\]
保留与未知参数有关的部分，即
\[
L(x\mid\sigma^2)\propto(\sigma^2)^{-n/2}\exp\left(-\frac{T}{2\sigma^2}\right).
\]
\step{Step 2：先验密度}
取 $\sigma^2\sim\operatorname{IGa}(\alpha,\lambda)$，$\alpha,\lambda>0$，则
\[
\pi(\sigma^2)=\frac{\lambda^\alpha}{\Gamma(\alpha)}
(\sigma^2)^{-\alpha-1}e^{-\lambda/\sigma^2},\qquad\sigma^2>0.
\]
\step{Step 3：相乘、整理幂次与指数}
\begin{align*}
\pi(\sigma^2\mid x)&\propto
(\sigma^2)^{-n/2}e^{-T/(2\sigma^2)}
(\sigma^2)^{-\alpha-1}e^{-\lambda/\sigma^2}\\
&=(\sigma^2)^{-(\alpha+n/2)-1}
\exp\left[-\frac{\lambda+T/2}{\sigma^2}\right].
\end{align*}
\step{Step 4：识别并归一化后验}
令 $a=\alpha+n/2$，$b=\lambda+T/2$，则归一化后的后验密度为
\[
\pi(\sigma^2\mid x)=\frac{b^a}{\Gamma(a)}
(\sigma^2)^{-a-1}e^{-b/\sigma^2},\quad\sigma^2>0.
\]
因此
\[
\sigma^2\mid x\sim\operatorname{IGa}\left(\alpha+\frac n2,\lambda+\frac T2\right).
\]
\step{Step 5：推导后验均值}
为避免把方差参数与标准差混淆，令 $v=\sigma^2$ 为积分变量。当 $a>1$ 时，代换 $u=b/v$，有
\begin{align*}
E(\sigma^2\mid x)
&=\frac{b^a}{\Gamma(a)}\int_0^\infty v^{-a}e^{-b/v}\,dv\\
&=\frac{b^a}{\Gamma(a)}b^{1-a}\int_0^\infty u^{a-2}e^{-u}\,du\\
&=b\frac{\Gamma(a-1)}{\Gamma(a)}=\frac b{a-1}.
\end{align*}
所以当 $\alpha+n/2>1$ 时，后验均值为
\[
E(\sigma^2\mid x)=\frac{\lambda+T/2}{\alpha+n/2-1}.
\]
当 $\alpha>1$ 时，先验均值为 $\lambda/(\alpha-1)$，并可改写为
\[
E(\sigma^2\mid x)=\gamma\frac{\lambda}{\alpha-1}+(1-\gamma)\frac Tn,
\qquad \gamma=\frac{\alpha-1}{\alpha+n/2-1}.
\]
这个加权式由分子拆开得到。记 $D=\alpha-1+n/2$，则
\[
\frac{\lambda+T/2}{D}
=\frac{\alpha-1}{D}\frac{\lambda}{\alpha-1}
+\frac{n/2}{D}\frac Tn,
\qquad\frac{\alpha-1}{D}+\frac{n/2}{D}=1.
\]
这里 $T/n$ 使用已知的总体均值 $\mu$ 计算；它不是以 $\bar x$ 为中心、分母为 $n-1$ 的样本方差。上述先验均值的加权解释要求 $\alpha>1$；后验均值本身只要求 $\alpha+n/2>1$。

\section{超参数及其确定}
先验分布中未知的参数称为超参数。例如 $\operatorname{Be}(\alpha,\beta)$ 中 $\alpha,\beta$，以及 $N(\mu,\tau^2)$ 中的 $\mu,\tau^2$。

\subsection{利用先验矩}
若历史信息给出 $\theta$ 的先验均值 $\bar\theta$ 与方差 $s_\theta^2$，并用 $\operatorname{Be}(\alpha,\beta)$ 表示先验，则令
\[
\bar\theta=\frac{\alpha}{\alpha+\beta},\qquad
s_\theta^2=\frac{\alpha\beta}{(\alpha+\beta)^2(\alpha+\beta+1)}.
\]
解得
\[
\alpha=\bar\theta\left(\frac{(1-\bar\theta)\bar\theta}{s_\theta^2}-1\right),\qquad
\beta=(1-\bar\theta)\left(\frac{(1-\bar\theta)\bar\theta}{s_\theta^2}-1\right).
\]

\subsection{利用分位数或混合信息}
若已知上下四分位数 $\theta_L,\theta_U$，可令贝塔分布满足
\[
P(\theta\leq\theta_L)=0.25,\qquad P(\theta\geq\theta_U)=0.25,
\]
数值求解 $\alpha,\beta$。若已知均值 $\bar\theta$ 和一个 $p$ 分位数 $\theta_p$，则解
\[
\frac{\alpha}{\alpha+\beta}=\bar\theta,\qquad
\int_0^{\theta_p}\frac{t^{\alpha-1}(1-t)^{\beta-1}}{B(\alpha,\beta)}dt=p.
\]
若只知道均值，则只能确定 $\alpha/(\alpha+\beta)$，不能唯一确定两个超参数。还需表达对均值的确信程度；在均值固定时，$\alpha+\beta$ 越大，先验方差越小、先验越集中。

\section{多参数模型}
\subsection{正态均值与方差联合先验}
课堂涉及正态模型 $X_i\mid\mu,\sigma^2\sim N(\mu,\sigma^2)$ 的联合共轭构造：先给定 $\sigma^2$，取 $\mu\mid\sigma^2$ 为正态分布；再对 $\sigma^2$ 取倒 Gamma 先验。联合先验与正态似然相乘后仍可整理为同一类正态--倒 Gamma 形式。后续计算通常通过似然核与先验核相乘、对 $\mu$ 配方，并整理出 $\sigma^2$ 的幂次和指数项。

更具体地，设
\[
\mu\mid\sigma^2\sim N\left(\mu_0,\frac{\sigma^2}{\kappa_0}\right),\qquad
\sigma^2\sim\operatorname{IGa}\left(\frac{\nu_0}{2},\frac{\nu_0\sigma_0^2}{2}\right).
\]
由 $\pi(\mu,\sigma^2)=\pi(\mu\mid\sigma^2)\pi(\sigma^2)$，联合先验核为
\[
\pi(\mu,\sigma^2)\propto(\sigma^2)^{-(\nu_0/2+3/2)}
\exp\left\{-\frac{\nu_0\sigma_0^2+\kappa_0(\mu-\mu_0)^2}{2\sigma^2}\right\}.
\]
这里条件正态密度中的 $(\sigma^2)^{-1/2}$ 也必须保留，因为 $\sigma^2$ 同样是未知参数；不能把它当成常数略去。

令 $\bar x=n^{-1}\sum_i x_i$，$s^2=(n-1)^{-1}\sum_i(x_i-\bar x)^2$，配方后得到
\[
\mu_n=\frac{\kappa_0\mu_0+n\bar x}{\kappa_0+n},\quad
\kappa_n=\kappa_0+n,\quad \nu_n=\nu_0+n,
\]
\[
\nu_n\sigma_n^2=\nu_0\sigma_0^2+(n-1)s^2+
\frac{n\kappa_0}{\kappa_0+n}(\bar x-\mu_0)^2.
\]
因此联合后验仍为正态--倒 Gamma。积分消去 $\sigma^2$ 后，$\mu$ 的边缘后验为 Student $t$ 分布，自由度为 $\nu_n$，位置为 $\mu_n$，尺度平方为 $\sigma_n^2/\kappa_n$。

\subsection{多元正态均值}
设 $Y_1,\ldots,Y_n\sim N_d(\mu,\Sigma)$，其中协方差矩阵 $\Sigma$ 已知，均值先验为 $\mu\sim N_d(\mu_0,\Lambda_0)$。记 $\bar y=n^{-1}\sum_iY_i$，则
\[
\mu\mid y\sim N_d(\mu_n,\Lambda_n),\quad
\Lambda_n^{-1}=\Lambda_0^{-1}+n\Sigma^{-1},\quad
\mu_n=\Lambda_n(\Lambda_0^{-1}\mu_0+n\Sigma^{-1}\bar y).
\]
后验精度矩阵是先验精度矩阵与样本均值精度矩阵之和，后验均值由两者按精度加权。

\section{充分统计量}
若样本关于参数的似然可分解为 $L(x\mid\theta)=g(T(x),\theta)h(x)$，其中样本对参数的依赖仅通过 $T(x)$，则 $T(x)$ 为充分统计量。本章示例中：二项模型的成功次数 $x$ 充分；正态均值且方差已知时样本均值 $\bar x$ 充分；正态均值已知、方差未知时 $T=\sum_i(x_i-\mu)^2$ 充分。充分统计量压缩数据而不丢失关于参数的似然信息。

经典定义：给定 $T(x)=t$ 后，样本 $x$ 的条件分布若不依赖 $\theta$，则 $T$ 对 $\theta$ 充分。因子分解定理给出等价判别：
\[
L(x\mid\theta)=g(T(x),\theta)h(x).
\]
从贝叶斯角度看，若 $T$ 充分，则对于使后验有定义的先验，$\pi(\theta\mid x)=\pi(\theta\mid T(x))$。这表示利用完整样本与仅利用充分统计量得到相同的后验。单独对某一个特殊先验验证相等，一般还不足以反过来证明充分性。

事实上，若 $p(x\mid\theta)=h(x)g(T(x),\theta)$，则
\[
\pi(\theta\mid x)
=\frac{h(x)g(T(x),\theta)\pi(\theta)}{h(x)\int g(T(x),u)\pi(u)\,du}
=\frac{g(T(x),\theta)\pi(\theta)}{\int g(T(x),u)\pi(u)\,du},
\]
右端只通过 $T(x)$ 依赖样本。

例如 $X_i\sim\operatorname{Exp}(\theta)$ 的联合似然为
\[
L(x\mid\theta)=\theta^n\exp\left(-\theta\sum_i x_i\right),
\]
故 $T=\sum_iX_i$ 充分。正态样本 $N(\mu,\sigma^2)$ 中，$\mu,\sigma^2$ 均未知时，$T=(\bar X,Q)$ 充分，其中 $Q=\sum_i(X_i-\bar X)^2$。


\subsection{用定义证明：指数样本的和充分}
设 $X_1,\ldots,X_n\mid\theta$ 独立同分布，密度为
\[
f(x_i\mid\theta)=\theta e^{-\theta x_i}\ind_{\{x_i>0\}},\qquad\theta>0.
\]
此处 $\theta$ 是\textbf{率参数}。令 $T=\sum_{i=1}^nX_i$，证明其充分性时，老师在板书中先作变量变换，再计算条件密度。

\step{第一步：把最后一个变量换成样本和}
作变换
\[
(X_1,\ldots,X_n)\longmapsto(X_1,\ldots,X_{n-1},T),
\qquad x_n=t-\sum_{i=1}^{n-1}x_i.
\]
逆变换的雅可比行列式绝对值为 $1$，于是
\[
f(x_1,\ldots,x_{n-1},t\mid\theta)
=\theta^ne^{-\theta t}\ind_{D_t},
\quad D_t=\left\{x_i>0,\ \sum_{i=1}^{n-1}x_i<t\right\}.
\]

\step{第二步：写出 $T$ 的密度}
由相同率参数的 Gamma 分布可加性，
\[
T\mid\theta\sim\Ga(n,\theta),\qquad
f_T(t\mid\theta)=\frac{\theta^n}{\Gamma(n)}t^{n-1}e^{-\theta t},\quad t>0.
\]

\step{第三步：计算给定 $T=t$ 时的条件密度}
\begin{align*}
f(x_1,\ldots,x_{n-1}\mid t,\theta)
&=\frac{f(x_1,\ldots,x_{n-1},t\mid\theta)}{f_T(t\mid\theta)}\\
&=\frac{(n-1)!}{t^{n-1}}\ind_{D_t}.
\end{align*}
右端不含 $\theta$。且在给定 $T=t$ 后，$X_n=t-\sum_{i<n}X_i$ 被前 $n-1$ 个变量确定，故整个样本在给定 $T$ 后的条件分布与 $\theta$ 无关。因此 $T$ 是 $\theta$ 的充分统计量。

\subsection{用因子分解证明：正态样本的两个充分统计量}
设 $X_i\mid\mu,\sigma^2\overset{iid}{\sim}N(\mu,\sigma^2)$，$\mu,\sigma^2$ 都未知。令
\[
\bar x=\frac1n\sum_i x_i,\qquad Q=\sum_i(x_i-\bar x)^2.
\]
先展开平方和：
\begin{align*}
\sum_i(x_i-\mu)^2
&=\sum_i\{(x_i-\bar x)+(\bar x-\mu)\}^2\\
&=Q+2(\bar x-\mu)\underbrace{\sum_i(x_i-\bar x)}_{=0}
+n(\bar x-\mu)^2\\
&=Q+n(\bar x-\mu)^2.
\end{align*}
因此似然为
\[
L(\mu,\sigma^2;x)=(2\pi\sigma^2)^{-n/2}
\exp\left\{-\frac{Q+n(\bar x-\mu)^2}{2\sigma^2}\right\}
=g((\bar x,Q),\mu,\sigma^2)\,h(x),
\]
其中 $h(x)=1$。由因子分解定理，$T=(\bar X,Q)$ 是 $(\mu,\sigma^2)$ 的充分统计量。

\subsection{例1.6.2：只使用充分统计量求正态后验}
设 $X_i\mid\theta\overset{iid}{\sim}N(\theta,1)$，先验 $\theta\sim N(0,\tau^2)$。因为 $\bar X$ 充分，可以直接用 $\bar X$ 的抽样分布。

\step{Step 1：似然函数}
\[
\bar X\mid\theta\sim N\left(\theta,\frac1n\right),\qquad
L(\theta;\bar x)=\sqrt{\frac{n}{2\pi}}
\exp\left\{-\frac n2(\theta-\bar x)^2\right\}.
\]
\step{Step 2：先验密度}
\[
\pi(\theta)=\frac1{\sqrt{2\pi}\tau}
\exp\left(-\frac{\theta^2}{2\tau^2}\right),\qquad\tau^2>0.
\]
\step{Step 3：相乘并配方}
\begin{align*}
\pi(\theta\mid\bar x)&\propto
\exp\left\{-\frac12\left[n(\theta-\bar x)^2+\tau^{-2}\theta^2\right]\right\}\\
&\propto\exp\left\{-\frac{n+\tau^{-2}}2
\left(\theta-\frac{n\bar x}{n+\tau^{-2}}\right)^2\right\}.
\end{align*}
\step{Step 4：识别后验分布并求均值}
\[
\boxed{\theta\mid x\sim N\left(
\frac{n\tau^2}{1+n\tau^2}\bar x,\frac{\tau^2}{1+n\tau^2}\right).}
\]
其后验均值为
\[
E(\theta\mid x)=\frac{n\tau^2}{1+n\tau^2}\bar x
+\frac1{1+n\tau^2}\cdot0.
\]
这里配方系数为 $n+\tau^{-2}$，并非 $n+\tau^2$；写成分式时可将分子、分母同时乘以 $\tau^2$ 核对。$\pi(\theta\mid x)=\pi(\theta\mid\bar x)$，所以不必再使用完整样本逐项相乘。

\subsection*{似然记号与条件分布记号的区别}
老师在板书中写的 $L(\theta;x)$ 或 $L(\theta\mid x)$，强调“观测 $x$ 固定后，似然是参数 $\theta$ 的函数”。课件写 $L(x\mid\theta)$，沿用的是抽样密度 $p(x\mid\theta)$ 的顺序。两种似然写法所指的数值函数相同：
\[
L(\theta;x):=p(x\mid\theta).
\]
但是 $\pi(\theta\mid x)$ 表示\textbf{后验密度}，必须再乘先验并归一化：
\[
\pi(\theta\mid x)=\frac{L(\theta;x)\pi(\theta)}{\int L(t;x)\pi(t)\,dt}.
\]
因此改变似然记号的排列顺序，不是把抽样分布变成了后验分布；均值未知与方差未知都遵守同一个公式。

\chapter{贝叶斯推断}

\section{条件观点}
在样本 $x$ 已经观察到以后，用 $\pi(\theta\mid x)$ 描述参数的不确定性，并在这个后验分布下进行估计、检验和预测。这就是本章贯穿始终的\textbf{条件观点}。

给定 $x$ 后，估计误差中的期望是对 $\theta\mid x$ 求的；频率方法的重复抽样误差则通常固定 $\theta$，对可能出现的样本 $X$ 求期望。两者的条件对象不同。

\section{估计}
\subsection{三种后验点估计}
给定后验分布后，板书使用以下三种点估计：
\begin{align*}
\hat\theta_{MD}&=\mathop{\arg\max}_{\theta}\pi(\theta\mid x)
&&\text{（后验众数，最大后验估计，MAP）},\\
P(\theta\leq\hat\theta_{Me}\mid x)&\geq\tfrac12,\quad
P(\theta\geq\hat\theta_{Me}\mid x)\geq\tfrac12
&&\text{（后验中位数）},\\
\hat\theta_E&=E(\theta\mid x)
&&\text{（后验均值）}.
\end{align*}
若后验连续且分布函数严格递增，中位数满足
$\int_{-\infty}^{\hat\theta_{Me}}\pi(\theta\mid x)\,d\theta=1/2$。
在正态后验中三者相同；一般分布中不一定相同。

损失函数决定如何从后验分布中选取估计。在似然与先验已经固定时，改变损失函数不会改变后验分布。例如平方损失下取后验均值，绝对损失下取后验中位数。给定行动 $a$ 的后验风险为
\[
\rho(a\mid x)=E[L(\theta,a)\mid x].
\]

\subsection{例2.2.1：正态均值及智商测量}
一般地，若 $X_i\mid\theta\overset{iid}{\sim}N(\theta,\sigma^2)$，$\sigma^2$ 已知，先验 $\theta\sim N(\mu,\tau^2)$，则
\[
\theta\mid x\sim N(\mu_1,\sigma_1^2),\quad
\sigma_1^2=\left(\frac n{\sigma^2}+\frac1{\tau^2}\right)^{-1},\quad
\mu_1=\sigma_1^2\left(\frac{n\bar x}{\sigma^2}+\frac\mu{\tau^2}\right).
\]
因此 $\hat\theta_{MD}=\hat\theta_{Me}=\hat\theta_E=\mu_1$。

板书中的智商例子为单次观测
\[
X\mid\theta\sim N(\theta,100),\qquad\theta\sim N(100,225),\qquad x=115.
\]
\step{Step 1：似然函数}
\[
L(\theta;x)=\frac1{\sqrt{2\pi\cdot100}}
\exp\left[-\frac{(x-\theta)^2}{200}\right].
\]
\step{Step 2：先验密度}
\[
\pi(\theta)=\frac1{\sqrt{2\pi\cdot225}}
\exp\left[-\frac{(\theta-100)^2}{450}\right].
\]
\step{Step 3：相乘、配方并识别后验}
\[
\pi(\theta\mid x)\propto
\exp\left\{-\frac12\left[\frac{(x-\theta)^2}{100}
+\frac{(\theta-100)^2}{225}\right]\right\}.
\]
后验精度和均值分别为
\[
\frac1{\sigma_1^2}=\frac1{100}+\frac1{225}=\frac{13}{900},\qquad
\mu_1=\frac{900}{13}\left(\frac{x}{100}+\frac{100}{225}\right)
=\frac{9x+400}{13}.
\]
归一化后为正态密度，故
\[
\theta\mid x=115\sim N\left(\frac{1435}{13},\frac{900}{13}\right)
\approx N(110.3846,69.2308).
\]
\step{Step 4：点估计与后验标准差}
\[
\boxed{\hat\theta_{MD}=\hat\theta_{Me}=\hat\theta_E=110.3846,\qquad
\sigma_1=\sqrt{900/13}\approx8.3205.}
\]
后验均值还可写为 $\frac9{13}x+\frac4{13}\cdot100$。正态分布第二个参数是\textbf{方差}，不能把 $69.2308$ 当成标准差。

\subsection{例2.2.2：二项分布与 Beta 先验}
设 $X\mid\theta\sim\Bin(n,\theta)$，$x$ 表示不合格品数，先验为 $\theta\sim\Be(\alpha,\beta)$，$\alpha,\beta>0$。
\step{Step 1：似然函数}
\[
L(\theta;x)=\binom nx\theta^x(1-\theta)^{n-x},\qquad 0<\theta<1.
\]
\step{Step 2：先验密度}
\[
\pi(\theta)=\frac1{B(\alpha,\beta)}
\theta^{\alpha-1}(1-\theta)^{\beta-1}\ind_{(0,1)}(\theta).
\]
\step{Step 3：后验核与归一化}
\[
\pi(\theta\mid x)\propto
\theta^{\alpha+x-1}(1-\theta)^{\beta+n-x-1}.
\]
令 $a=\alpha+x$，$b=\beta+n-x$。由于核在 $(0,1)$ 上的积分为 $B(a,b)$，有
\[
\boxed{\pi(\theta\mid x)=\frac{\theta^{a-1}(1-\theta)^{b-1}}{B(a,b)},\qquad
\theta\mid x\sim\Be(a,b).}
\]
若需边缘概率，则 $m(x)=\binom nx B(a,b)/B(\alpha,\beta)$。

\step{Step 4：求最大后验估计}
当 $a>1,b>1$ 时，令 $\ell(\theta)=\log\pi(\theta\mid x)$，有
\[
\ell'(\theta)=\frac{a-1}{\theta}-\frac{b-1}{1-\theta},\qquad
\ell''(\theta)=-\frac{a-1}{\theta^2}-\frac{b-1}{(1-\theta)^2}<0.
\]
由 $\ell'(\theta)=0$ 解得
\[
(a-1)(1-\theta)=(b-1)\theta
\quad\Longrightarrow\quad
\boxed{\hat\theta_{MD}=\frac{a-1}{a+b-2}
=\frac{\alpha+x-1}{\alpha+\beta+n-2}.}
\]
这个内部众数公式要求 $a,b>1$。若后验向某个端点单调上升，则按端点极限取众数；均匀密度没有唯一众数，$a,b<1$ 时则可能两端都有峰，不能直接套内部公式。

\step{Step 5：用 Beta 函数求后验均值}
\begin{align*}
\hat\theta_E
&=\frac1{B(a,b)}\int_0^1\theta^a(1-\theta)^{b-1}\,d\theta
=\frac{B(a+1,b)}{B(a,b)}\\
&=\frac{\Gamma(a+1)\Gamma(b)}{\Gamma(a+b+1)}
\frac{\Gamma(a+b)}{\Gamma(a)\Gamma(b)}
=\frac{a}{a+b}.
\end{align*}
其中使用 $\Gamma(a+1)=a\Gamma(a)$，故
\[
\boxed{\hat\theta_E=\frac{\alpha+x}{\alpha+\beta+n}.}
\]
记先验强度 $s=\alpha+\beta$、先验均值 $\mu_0=\alpha/s$，则在 $n>0$ 时
\[
\hat\theta_E=\frac{s}{s+n}\mu_0+\frac n{s+n}\frac xn.
\]
即后验均值是先验均值与样本不合格品率的加权平均。

\step{课堂补充：未检出不合格品时的尾概率}
若先验均匀，即 $\alpha=\beta=1$，且观察到 $x=0$，则
\[
\theta\mid x=0\sim\Be(1,n+1),\qquad
\pi(\theta\mid x=0)=(n+1)(1-\theta)^n.
\]
对 $0<t<1$，
\[
P(\theta>t\mid x=0)=\int_t^1(n+1)(1-\theta)^n\,d\theta=(1-t)^{n+1}.
\]
取 $t=0.2$：$n=3$ 时为 $0.8^4=0.4096$；$n=10$ 时为 $0.8^{11}\approx0.08590$。同样“没有检出不合格品”，样本量不同，后验判断不同。

\subsection{例2.2.3：平移指数分布与柯西先验}
本例在板书上也记作“例2.4”。设 $X=\theta+Y$，$Y\sim\operatorname{Exp}(1)$。
\step{Step 1：似然函数及参数范围}
\[
L(\theta;x)=e^{-(x-\theta)}\ind_{\{\theta\leq x\}}.
\]
观测 $x$ 固定后，$\theta\leq x$ 是似然对参数的限制，必须保留。
\step{Step 2：先验密度}
\[
\pi(\theta)=\frac1{\pi(1+\theta^2)},\qquad-\infty<\theta<\infty.
\]
\step{Step 3：相乘并归一化}
\[
m(x)=\frac{e^{-x}}\pi\int_{-\infty}^x\frac{e^t}{1+t^2}\,dt,
\]
\[
\pi(\theta\mid x)=
\frac{e^\theta/(1+\theta^2)}{\displaystyle\int_{-\infty}^x e^t/(1+t^2)\,dt}
\ind_{\{\theta\leq x\}}.
\]
\step{Step 4：求众数与均值}
令 $g(\theta)=e^\theta/(1+\theta^2)$，则
\[
g'(\theta)=\frac{e^\theta(1+\theta^2)-2\theta e^\theta}{(1+\theta^2)^2}
=\frac{e^\theta(\theta-1)^2}{(1+\theta^2)^2}\geq0.
\]
因此后验密度在允许范围内单调不减，其最大值在 $\theta=x$ 处取得：
\[
\boxed{\hat\theta_{MD}=x=\hat\theta_{MLE}.}
\]
但后验均值为
\[
\hat\theta_E=
\frac{\displaystyle\int_{-\infty}^x\frac{t e^t}{1+t^2}\,dt}
{\displaystyle\int_{-\infty}^x\frac{e^t}{1+t^2}\,dt}<x.
\]
后验分布集中于 $(-\infty,x]$，且不是全都集中在 $x$，所以其均值严格小于 $x$；不能因 MAP 等于 $x$ 就把后验均值也写成 $x$。

\subsection{后验估计误差：板书中的完整展开}
固定已经观察到的 $x$，估计值 $\hat\theta=\hat\theta(x)$ 是常数。后验均方误差定义为
\[
\MSE(\hat\theta\mid x)=E[(\theta-\hat\theta)^2\mid x]
=\int(\theta-\hat\theta)^2\pi(\theta\mid x)\,d\theta.
\]
令 $m=E(\theta\mid x)$，则
\[
\theta-\hat\theta=(\theta-m)+(m-\hat\theta).
\]
平方并求后验期望：
\begin{align*}
\MSE(\hat\theta\mid x)
&=E[(\theta-m)^2\mid x]
+2(m-\hat\theta)E(\theta-m\mid x)+(m-\hat\theta)^2\\
&=\Var(\theta\mid x)+0+(m-\hat\theta)^2.
\end{align*}
所以
\[
\boxed{\MSE(\hat\theta\mid x)=\Var(\theta\mid x)
+\{E(\theta\mid x)-\hat\theta\}^2.}
\]
这里交叉项等于零，是因为 $E(\theta-m\mid x)=m-m=0$。于是
\[
\MSE(\hat\theta\mid x)\geq\Var(\theta\mid x),
\]
等号在 $\hat\theta=\hat\theta_E=m$ 时取得。这就证明了\textbf{后验均值在平方损失下最优}。

估计值 $\hat\theta$ 的后验标准误为 $\sqrt{\MSE(\hat\theta\mid x)}$；后验标准差为 $\sqrt{\Var(\theta\mid x)}$。取后验均值作为估计时二者相等，一般情况下不相等。

\subsection{例2.2.4：几何抽检与离散后验}
本例在板书上也记作“例2.5”。逐件检查产品，直到发现第一个不合格品；$X$ 是累计检查件数，$\theta$ 是不合格品率。已知
\[
\theta\in\left\{\tfrac14,\tfrac12,\tfrac34\right\},\qquad
P(\theta=\tfrac14)=P(\theta=\tfrac12)=P(\theta=\tfrac34)=\tfrac13,
\]
观察到 $X=3$。
\step{Step 1：似然}
\[
P(X=x\mid\theta)=\theta(1-\theta)^{x-1},\qquad x=1,2,\ldots.
\]
故 $L(\theta;3)=\theta(1-\theta)^2$。
\step{Step 2：先验与联合概率}
先验三点各为 $1/3$，似然乘先验得到
\[
\begin{array}{c|ccc}
\theta&1/4&1/2&3/4\\\hline
P(X=3\mid\theta)&9/64&1/8&3/64\\
P(X=3,\theta)&9/192&8/192&3/192
\end{array}
\]
\step{Step 3：边缘概率及归一化}
\[
m(3)=\frac{9+8+3}{192}=\frac5{48},\qquad
P(\theta=t\mid X=3)=\frac{P(X=3,\theta=t)}{m(3)}.
\]
所以
\[
\begin{array}{c|ccc}
\theta&1/4&1/2&3/4\\\hline
P(\theta\mid X=3)&9/20&8/20&3/20
\end{array}
\]
\step{Step 4：点估计、后验矩与误差}
\[
\hat\theta_{MD}=\frac14,\qquad
\hat\theta_E=\frac14\frac9{20}+\frac12\frac8{20}+\frac34\frac3{20}
=\frac{17}{40}.
\]
累积到 $1/4$ 时为 $9/20<1/2$，累积到 $1/2$ 时为 $17/20>1/2$，故后验中位数为 $\hat\theta_{Me}=1/2$。
\[
E(\theta^2\mid3)=\frac1{16}\frac9{20}+\frac14\frac8{20}+\frac9{16}\frac3{20}
=\frac{17}{80},
\]
\[
\Var(\theta\mid3)=\frac{17}{80}-\left(\frac{17}{40}\right)^2
=\frac{51}{1600}.
\]
因此
\[
\MSE(\hat\theta_E\mid3)=\frac{51}{1600},\qquad
\MSE(\hat\theta_{MD}\mid3)=\frac{51}{1600}
+\left(\frac{17}{40}-\frac14\right)^2=\frac1{16}.
\]
后验众数估计的后验标准误为 $\sqrt{1/16}=1/4$。也可以直接计算
\[
E[(\theta-\tfrac14)^2\mid3]
=0^2\frac9{20}+(\tfrac14)^2\frac8{20}+(\tfrac12)^2\frac3{20}
=\frac1{16},
\]
与分解公式所得结果相同。

\subsection{例2.2.5：Beta 后验下两种估计的误差比较}
在二项模型中，取均匀先验 $\Be(1,1)$，则
\[
\theta\mid x\sim\Be(x+1,n-x+1),\qquad
\hat\theta_E=\frac{x+1}{n+2},\quad\hat\theta_{MD}=\frac xn\quad(n>0).
\]
其中 $x=0$ 或 $n$ 时，众数按端点极限解释。后验方差为
\[
V_x:=\Var(\theta\mid x)=\frac{(x+1)(n-x+1)}{(n+2)^2(n+3)}.
\]
因此
\[
\MSE(\hat\theta_E\mid x)=V_x,\qquad
\MSE(\hat\theta_{MD}\mid x)=V_x+
\left(\frac{x+1}{n+2}-\frac xn\right)^2.
\]
通分得
\[
\frac{x+1}{n+2}-\frac xn=\frac{n-2x}{n(n+2)},
\]
所以
\[
\boxed{\MSE(\hat\theta_{MD}\mid x)-\MSE(\hat\theta_E\mid x)
=\frac{(n-2x)^2}{n^2(n+2)^2}\geq0.}
\]
课堂代入 $n=3,x=0$：
\[
\hat\theta_E=\frac15,\qquad\hat\theta_{MD}=0,\qquad
\Var(\theta\mid x)=\frac{1\cdot4}{5^2\cdot6}=\frac2{75},
\]
\[
\MSE(\hat\theta_{MD}\mid x)=\frac2{75}+(\tfrac15-0)^2=\frac1{15}.
\]

\subsection{例2.2.6：均匀分布上限参数的后验均值（课件补全）}
设 $X_1,\ldots,X_n\mid\theta\overset{iid}{\sim}U(0,\theta)$，先验 $\theta\sim\IGa(\alpha,\lambda)$，$\alpha,\lambda>0$。求 $E(\theta\mid x)$。
\step{Step 1：直接从完整样本写似然}
假设观测均为正数，记 $y=\max_i x_i$，则
\[
L(\theta;x)=\prod_{i=1}^n\frac1\theta\ind_{\{0<x_i<\theta\}}
=\theta^{-n}\ind_{\{\theta>y\}}.
\]
这里\textbf{不必先求最大次序统计量的分布}。引入 $y$ 只是简写“$\theta$ 必须大于所有观测值”；也可一直写 $\max_i x_i$。课件用 $Y=\max_iX_i$ 的密度 $ny^{n-1}\theta^{-n}\ind_{\{0<y<\theta\}}$ 求解，其多出的 $ny^{n-1}$ 与 $\theta$ 无关，所得后验相同。
\step{Step 2：先验密度}
\[
\pi(\theta)=\frac{\lambda^\alpha}{\Gamma(\alpha)}
\theta^{-\alpha-1}e^{-\lambda/\theta}\ind_{\{\theta>0\}}.
\]
\step{Step 3：相乘并保留后验的支撑范围}
令 $a=\alpha+n$，则
\[
\pi(\theta\mid x)\propto\theta^{-a-1}e^{-\lambda/\theta}\ind_{\{\theta>y\}}.
\]
它是 $\IGa(a,\lambda)$ \textbf{在 $(y,\infty)$ 上截断后的分布}，不是整个 $(0,\infty)$ 上的普通倒 Gamma 分布。
\step{Step 4：归一化及均值的积分比}
定义下不完全 Gamma 函数 $\gamma(s,z)=\int_0^z u^{s-1}e^{-u}\,du$。代换 $u=\lambda/\theta$，得
\[
D=\int_y^\infty\theta^{-a-1}e^{-\lambda/\theta}\,d\theta
=\lambda^{-a}\gamma(a,\lambda/y).
\]
所以后验密度为
\[
\pi(\theta\mid x)=\frac{\lambda^a}{\gamma(a,\lambda/y)}
\theta^{-a-1}e^{-\lambda/\theta}\ind_{\{\theta>y\}}.
\]
当 $a>1$ 时，
\begin{align*}
E(\theta\mid x)
&=\frac{\displaystyle\int_y^\infty\theta^{-a}e^{-\lambda/\theta}\,d\theta}
{\displaystyle\int_y^\infty\theta^{-a-1}e^{-\lambda/\theta}\,d\theta}\\
&=\lambda\frac{\gamma(a-1,\lambda/y)}{\gamma(a,\lambda/y)}.
\end{align*}
\step{Step 5：按课件化为卡方分布函数之比}
因为
\[
P(\chi^2_{2s}\leq v)=\frac{\gamma(s,v/2)}{\Gamma(s)},\qquad
\Gamma(a)=(a-1)\Gamma(a-1),
\]
故
\[
\boxed{E(\theta\mid x)=\frac{\lambda}{\alpha+n-1}
\frac{P(\chi^2_{2(\alpha+n-1)}\leq2\lambda/y)}
{P(\chi^2_{2(\alpha+n)}\leq2\lambda/y)}.}
\]
卡方自由度是 $2(\alpha+n-1)$ 与 $2(\alpha+n)$，不能漏掉因子 $2$。

课件数值 $n=5,y=2,\alpha=2,\lambda=8$ 给出 $2\lambda/y=8$：
\[
E(\theta\mid x)=\frac86\frac{P(\chi^2_{12}\leq8)}{P(\chi^2_{14}\leq8)}
=\frac86\frac{0.2148696\cdots}{0.1106740\cdots}\approx2.5886.
\]
课件采用舍入概率给出约 $2.5884$，差别来自数值舍入。计算时必须保留截断所带来的两个分布函数。

\section{区间估计}
\subsection{可信区间、分位数与单侧可信限}
给定观测 $x$，若区间 $[L(x),U(x)]$ 满足
\[
P\{L(x)\leq\theta\leq U(x)\mid x\}\geq1-\alpha,
\]
则称其为可信水平 $1-\alpha$ 的可信区间。若后验连续，上式左端写为
\[
\int_{L(x)}^{U(x)}\pi(\theta\mid x)\,d\theta.
\]
板书保留“$\geq$”是为了兼顾离散分布：离散概率不一定能恰好凑成 $1-\alpha$。例如
\[
P(\theta=0\mid x)=0.2,\quad P(\theta=1\mid x)=0.5,\quad P(\theta=2\mid x)=0.3,
\]
此时 $[0,2]$ 的后验概率为 $1\geq0.9$，可作 $90\%$ 可信区间。

记后验分布函数和分位数为
\[
F_x(t)=P(\theta\leq t\mid x),\qquad
q_p(x)=\inf\{t:F_x(t)\geq p\}.
\]
\textbf{等尾可信区间}取
\[
\boxed{[q_{\alpha/2}(x),q_{1-\alpha/2}(x)].}
\]
当后验连续且相应分位数处无平坦段时，两侧尾概率各为 $\alpha/2$。

单侧情形中，若
\[
P(\theta\geq L\mid x)\geq1-\alpha,
\]
则 $L$ 为可信水平 $1-\alpha$ 的\textbf{单侧下可信限}；若
\[
P(\theta\leq U\mid x)\geq1-\alpha,
\]
则 $U$ 为\textbf{单侧上可信限}。连续且严格递增的分布中可分别取 $L=q_\alpha$ 与 $U=q_{1-\alpha}$。

\subsection{例2.3.1：正态后验的可信区间}
若 $\theta\mid x\sim N(\mu_1,\sigma_1^2)$，则
\[
\frac{\theta-\mu_1}{\sigma_1}\biggm|x\sim N(0,1).
\]
以下用 $z_p=\Phi^{-1}(p)$ 表示标准正态的\textbf{下侧 $p$ 分位数}，即 $P(Z\leq z_p)=p$。从而
\[
P\left\{-z_{1-\alpha/2}\leq\frac{\theta-\mu_1}{\sigma_1}
\leq z_{1-\alpha/2}\mid x\right\}=1-\alpha,
\]
得到
\[
\boxed{[\mu_1-z_{1-\alpha/2}\sigma_1,\ \mu_1+z_{1-\alpha/2}\sigma_1].}
\]
若教材用 $z_{\alpha/2}$ 表示\textbf{上尾概率}为 $\alpha/2$ 的正数，则数值相同，只是分位数记法不同。

对智商例子，$\mu_1=110.3846$，$\sigma_1=8.3205$，故 $95\%$ 可信区间为
\[
110.3846\pm1.96\times8.3205\approx[94.08,126.69].
\]
这是给定本次测量后关于真实参数 $\theta$ 的区间，不是下一次测量值的预测区间。

\subsection{例2.3.2：截尾寿命试验的倒 Gamma 后验}
电视机寿命 $T$ 在给定平均寿命 $\theta$ 下服从指数分布。本例 $\theta$ 是\textbf{均值参数}，与第一章指数分布的率参数记法应区分。

\step{Step 1：似然函数}
寿命密度与存活概率为
\[
p(t\mid\theta)=\frac1\theta e^{-t/\theta},\quad t>0,
\qquad P(T>t\mid\theta)=e^{-t/\theta}.
\]
在 $n$ 台电视机中，记录到 $r$ 台在时刻 $t_1\leq\cdots\leq t_r$ 失效，其余 $n-r$ 台在停止时刻 $c$ 仍存活，故
\begin{align*}
L(\theta;t)
&\propto\prod_{i=1}^r\left(\frac1\theta e^{-t_i/\theta}\right)
\left(e^{-c/\theta}\right)^{n-r}\\
&=\theta^{-r}\exp\left\{-\frac{\sum_{i=1}^r t_i+(n-r)c}{\theta}\right\}.
\end{align*}
令总试验时间 $S_r=\sum_{i=1}^r t_i+(n-r)c$。定数截尾时 $c=t_r$，这就得到板书的
\[
S_r=\sum_{i=1}^r t_i+(n-r)t_r.
\]
若一次失效也没有（$r=0$），应理解为在预定时刻 $c$ 停止，此时 $S_0=nc$；不能写不存在的 $t_0$。

\step{Step 2：先验密度}
取 $\theta\sim\IGa(\alpha,\beta)$：
\[
\pi(\theta)=\frac{\beta^\alpha}{\Gamma(\alpha)}
\theta^{-\alpha-1}e^{-\beta/\theta},\quad\theta>0.
\]
\step{Step 3：相乘并识别后验}
\[
\pi(\theta\mid t)\propto\theta^{-r}e^{-S_r/\theta}
\theta^{-\alpha-1}e^{-\beta/\theta}
=\theta^{-(\alpha+r)-1}e^{-(\beta+S_r)/\theta}.
\]
因此
\[
\boxed{\theta\mid t\sim\IGa(\alpha+r,\beta+S_r).}
\]
其归一化常数是 $(\beta+S_r)^{\alpha+r}/\Gamma(\alpha+r)$。
\step{Step 4：后验均值及课堂数值}
当 $\alpha+r>1$ 时，
\[
\hat\theta_E=\frac{\beta+S_r}{\alpha+r-1}.
\]
沿用板书 $r=0,S_0=40000,\alpha=1.956,\beta=2868$，则
\[
\theta\mid t\sim\IGa(1.956,42868),\qquad
\hat\theta_E=\frac{42868}{0.956}\approx44841\text{小时}.
\]
\textit{数值核对：}这里按板书取 $\beta=2868$。课件先验确定部分另写“均值为30000小时”，但 $2868/(1.956-1)=3000$，两者相差一位；不能把这两项当成同时成立的条件。本例后续计算沿用板书参数及后验 $42868$。

\subsection{最大后验密度可信集与 HPD 区间}
可信水平相同的区间可以有许多个。老师进一步问：能否优先选取后验密度较高的部分，使区间尽可能短？

记 $f(\theta)=\pi(\theta\mid x)$，对阈值 $k$ 定义
\[
C_k=\{\theta:f(\theta)\geq k\}.
\]
选择 $k$ 使
\[
\int_{C_k}f(\theta)\,d\theta=1-\alpha,
\]
所得 $C_k$ 称最大后验密度（HPD）可信集。\textbf{单峰、连续且无平顶等常见条件下}，$C_k$ 是区间；若后验多峰，HPD 集可能是若干段，不一定是单个区间。

\step{板书推导：为什么两个内部端点的密度相等？}
在单峰情形中，设区间为 $[L,U]$，有
\[
F_x(U)-F_x(L)=1-\alpha.
\]
在保持后验概率不变的条件下，把 $U$ 看成 $L$ 的函数，求导得
\[
f(U)\frac{dU}{dL}-f(L)=0,
\qquad \frac{dU}{dL}=\frac{f(L)}{f(U)}.
\]
区间长度为 $U-L$，内部最优点满足
\[
\frac d{dL}(U-L)=\frac{f(L)}{f(U)}-1=0.
\]
故需联立
\[
\boxed{\begin{cases}
\displaystyle\int_L^U f(\theta)\,d\theta=1-\alpha,\\
f(L)=f(U).
\end{cases}}
\]
若最优端点位于支撑边界，上述内部端点等密度条件不一定适用。

\step{板书中的阈值搜索法}
定义 $M(k)=\int_{\{f\geq k\}}f(\theta)\,d\theta$。若 $M(k)>1-\alpha$，表示取入的概率过多，应增大 $k$；若 $M(k)<1-\alpha$，则减小 $k$。每次由 $f(L)=f(U)=k$ 找到两端点，迭代直到概率满足要求。

\step{寿命例子的 $90\%$ HPD 区间}
对 $\IGa(a,b)$，对数密度中与 $\theta$ 有关的部分为
\[
-(a+1)\log\theta-\frac b\theta.
\]
令 $h(\theta)=(a+1)\log\theta+b/\theta$，则 $f(L)=f(U)$ 等价于 $h(L)=h(U)$，数值上更便于计算。众数 $b/(a+1)$ 位于两端点之间。

\textbf{为了对应老师板书的数值，这一步按课件将 $a=1.956$ 近似为 $2$，保持 $b=42868$。}此时
\[
f(\theta)=b^2\theta^{-3}e^{-b/\theta},\qquad
F(\theta)=\left(1+\frac b\theta\right)e^{-b/\theta}.
\]
计算程序于是化为求解
\[
\begin{cases}
\displaystyle\left(1+\frac bU\right)e^{-b/U}
-\left(1+\frac bL\right)e^{-b/L}=0.9,\\
3\log L+b/L=3\log U+b/U.
\end{cases}
\]
所得
\[
\boxed{[L,U]\approx[4735.17,81183.22]\text{小时}\quad(a\approx2).}
\]
保留原后验 $\IGa(1.956,42868)$ 而不作形状参数近似时，数值解约为 $[4750.48,84767.95]$ 小时。两个区间对应不同的形状参数，不能混用。

\section{假设检验}
\subsection{后验概率、边缘似然与检验决策}
考虑
\[
H_0:\theta\in\Theta_0,\qquad H_1:\theta\in\Theta_1,
\quad\Theta_0\cap\Theta_1=\varnothing,\quad\Theta_0\cup\Theta_1=\Theta.
\]
按板书记假设的先验概率和后验概率分别为
\[
\pi_i=P(H_i)>0,\qquad \alpha_i=P(H_i\mid x),\qquad i=0,1.
\]
这里 $\alpha_i$ 是\textbf{假设的后验概率}，不是可信水平中的尾概率 $\alpha$，也不是频率检验中的显著性水平。

若 $\pi(\theta)$ 是连续先验，则在假设 $H_i$ 成立的条件下，参数的条件先验密度为
\[
g_i(\theta)=\pi(\theta\mid H_i)
=\frac{\pi(\theta)\ind_{\Theta_i}(\theta)}{\pi_i}.
\]
总体先验写为 $\pi(\theta)=\pi_0g_0(\theta)+\pi_1g_1(\theta)$。定义 $H_i$ 下数据的边缘密度为
\[
m_i(x)=\int_{\Theta_i}p(x\mid\theta)g_i(\theta)\,d\theta.
\]
于是由事件形式贝叶斯公式，
\[
\boxed{\alpha_i=\frac{\pi_i m_i(x)}{\pi_0m_0(x)+\pi_1m_1(x)}.}
\]
若已求出完整后验，也可以直接积分：
\[
\alpha_i=\int_{\Theta_i}\pi(\theta\mid x)\,d\theta,\qquad
\alpha_0+\alpha_1=1.
\]

\step{两种错误损失相同时}
正确决策损失为零，两种错误损失都取 $1$。选择 $H_0$ 的后验风险为 $\alpha_1$，选择 $H_1$ 的后验风险为 $\alpha_0$。因此
\[
\alpha_0>\alpha_1\Rightarrow\text{选择 }H_0,\qquad
\alpha_1>\alpha_0\Rightarrow\text{选择 }H_1.
\]
相等时两种决策的后验风险相同。

\step{错误损失不同时：板书反例}
令 $C_{01}$ 表示“选择 $H_0$、实际 $H_1$ 为真”的损失，$C_{10}$ 表示反向错误的损失，正确决策仍取零。则
\[
\rho_0(x)=C_{01}\alpha_1,\qquad\rho_1(x)=C_{10}\alpha_0.
\]
若 $\alpha_0=0.8,\alpha_1=0.2$，但 $C_{01}=10,C_{10}=1$，则
\[
\rho_0=10\times0.2=2,\qquad\rho_1=1\times0.8=0.8,
\]
故应选择 $H_1$。不能脱离损失条件，机械地“选择后验概率较大的假设”。

\step{与频率检验的区别}
板书另写频率方法对拒绝域 $W$ 的约束
\[
\sup_{\theta\in\Theta_0}P_\theta(X\in W)\leq\alpha.
\]
这是在 $H_0$ 下控制重复抽样误拒概率；后验概率 $P(H_0\mid x)$ 是给定已经观察到的数据后关于假设的概率，两者含义不同。

\subsection{贝叶斯因子：后验赔率除以先验赔率}
由后验概率公式相除，
\[
\frac{\alpha_0}{\alpha_1}
=\frac{\pi_0}{\pi_1}\frac{m_0(x)}{m_1(x)}.
\]
支持 $H_0$ 相对于 $H_1$ 的贝叶斯因子定义为
\[
\boxed{\BF_{01}(x)=\frac{\alpha_0/\alpha_1}{\pi_0/\pi_1}
=\frac{m_0(x)}{m_1(x)}.}
\]
即
\[
\boxed{\text{后验赔率}=\text{先验赔率}\times\text{贝叶斯因子}.}
\]
$\BF_{01}>1$ 表示数据使相对证据向 $H_0$ 移动，$\BF_{01}<1$ 则向 $H_1$ 移动，且 $\BF_{10}=1/\BF_{01}$。\textbf{贝叶斯因子大于1本身不保证最终选择 $H_0$}，还需结合先验赔率和损失。

当先验概率相同、两种错误损失相同时，可以按 $\BF_{01}$ 与 $1$ 比较。一般地，选择 $H_0$ 的条件为
\[
\frac{\alpha_0}{\alpha_1}>\frac{C_{01}}{C_{10}}
\quad\Longleftrightarrow\quad
\BF_{01}>\frac{C_{01}}{C_{10}}\frac{\pi_1}{\pi_0}.
\]

\subsection{例2.4.1：二项概率的单侧检验与三种积分算法}
板书也记作“例2.10”。设 $X\mid\theta\sim\Bin(n,\theta)$，先验为 $\theta\sim U(0,1)$，检验
\[
H_0:\theta\leq\tfrac12,\qquad H_1:\theta>\tfrac12.
\]
\step{Step 1：似然函数}
\[
L(\theta;x)=\binom nx\theta^x(1-\theta)^{n-x}.
\]
\step{Step 2：先验及两个假设的先验概率}
先验密度为 $\pi(\theta)=1$，$0<\theta<1$，所以
\[
\pi_0=\pi_1=\frac12,
\]
在两半区间上的条件先验密度分别为 $2\ind_{(0,1/2]}$ 与 $2\ind_{(1/2,1)}$。
\step{Step 3：后验分布}
\[
\pi(\theta\mid x)\propto\theta^x(1-\theta)^{n-x},\qquad
\theta\mid x\sim\Be(x+1,n-x+1),
\]
\[
\pi(\theta\mid x)=\frac{(n+1)!}{x!(n-x)!}\theta^x(1-\theta)^{n-x}.
\]
\step{Step 4：计算后验假设概率}
\[
\alpha_0=\frac{(n+1)!}{x!(n-x)!}
\int_0^{1/2}\theta^x(1-\theta)^{n-x}\,d\theta,
\qquad\alpha_1=1-\alpha_0.
\]
老师在板书中给出了下面三种算法。

\step{方法一：分部积分，建立递推}
记 $I_{p,q}=\int_0^{1/2}\theta^p(1-\theta)^q\,d\theta$，取
$u=(1-\theta)^q$，$dv=\theta^p\,d\theta$，得到
\begin{align*}
I_{p,q}
&=\left[\frac{\theta^{p+1}(1-\theta)^q}{p+1}\right]_0^{1/2}
+\frac q{p+1}\int_0^{1/2}\theta^{p+1}(1-\theta)^{q-1}\,d\theta\\
&=\frac{2^{-(p+q+1)}}{p+1}+\frac q{p+1}I_{p+1,q-1}.
\end{align*}
一直递推到 $I_{p+q,0}=2^{-(p+q+1)}/(p+q+1)$ 即可。
代入课堂数值 $n=5,x=1$：
\begin{align*}
I_{1,4}
&=\frac1{128}\left(1+\frac43+\frac{4\cdot3}{3\cdot4}
+\frac{4\cdot3\cdot2}{3\cdot4\cdot5}
+\frac{4!}{3\cdot4\cdot5\cdot6}\right)\\
&=\frac1{128}\left(1+\frac43+1+\frac25+\frac1{15}\right)
=\frac{19}{640}.
\end{align*}
因此
\[
\alpha_0=\frac{6!}{1!4!}\frac{19}{640}=\frac{57}{64},\qquad
\alpha_1=\frac7{64}.
\]

\step{方法二：二项展开后逐项积分}
令 $m=n-x$，由 $(1-\theta)^m=\sum_{k=0}^m(-1)^k\binom mk\theta^k$，
\[
I_{x,m}=\sum_{k=0}^m(-1)^k\binom mk
\frac{2^{-(x+k+1)}}{x+k+1}.
\]
对 $x=1,m=4$，
\[
I_{1,4}=\frac18-\frac16+\frac3{32}-\frac1{40}+\frac1{384}
=\frac{57}{1920}=\frac{19}{640},
\]
再乘 $30$ 即得 $\alpha_0=57/64$。

\step{方法三：Beta 分布函数与二项尾概率的关系}
当 $a,b$ 是正整数，$U\sim\Be(a,b)$ 时，
\[
P(U\leq t)=\sum_{j=a}^{a+b-1}\binom{a+b-1}{j}t^j(1-t)^{a+b-1-j}
=P\{\Bin(a+b-1,t)\geq a\}.
\]
代入 $a=x+1,b=n-x+1,t=1/2$，得
\[
\boxed{\alpha_0=2^{-(n+1)}\sum_{j=x+1}^{n+1}\binom{n+1}{j}.}
\]
例如 $n=5,x=1$，
\[
\alpha_0=P\{\Bin(6,\tfrac12)\geq2\}
=1-\frac{\binom60+\binom61}{64}=\frac{57}{64}.
\]

\step{Step 5：贝叶斯因子与决策}
先验赔率为 $1$，所以 $\BF_{01}=\alpha_0/\alpha_1$。当 $n=5$ 时，板书的对照表为
\begin{center}
\begin{tabular}{cccl}
\toprule
$x$&$P(H_0\mid x)$&$P(H_1\mid x)$&$\BF_{01}$\\
\midrule
0&$63/64$&$1/64$&$63$\\
1&$57/64$&$7/64$&$57/7\approx8.14$\\
2&$42/64$&$22/64$&$21/11\approx1.91$\\
3&$22/64$&$42/64$&$11/21\approx0.524$\\
4&$7/64$&$57/64$&$7/57\approx0.123$\\
5&$1/64$&$63/64$&$1/63\approx0.0159$\\
\bottomrule
\end{tabular}
\end{center}
在两种错误损失相同的条件下，$x=0,1,2$ 时选 $H_0$，$x=3,4,5$ 时选 $H_1$。还可利用对称关系
\[
\alpha_0(x)=\alpha_1(n-x)
\]
检查计算，减少重复。

\subsection{简单假设对简单假设：似然比形式}
若 $H_0:\theta=\theta_0$，$H_1:\theta=\theta_1$，则两个假设内没有待积分的参数，故
\[
m_0(x)=p(x\mid\theta_0),\qquad m_1(x)=p(x\mid\theta_1),\qquad
\BF_{01}(x)=\frac{p(x\mid\theta_0)}{p(x\mid\theta_1)}.
\]
两种错误损失相同时，选择 $H_1$ 的条件等价于
\[
\frac{p(x\mid\theta_1)}{p(x\mid\theta_0)}>\frac{\pi_0}{\pi_1}.
\]
这是一种似然比判别形式。若似然可分解为 $h(x)g(T(x),\theta)$，则 $h(x)$ 在比值中抵消，只需使用充分统计量。

\subsection{例2.4.2：正态均值的两个简单假设}
板书也记作“例2.11”。设 $X_i\mid\theta\overset{iid}{\sim}N(\theta,1)$，检验
\[
H_0:\theta=0,\qquad H_1:\theta=1.
\]
\step{Step 1：似然函数与平方和分解}
\[
L(\theta;x)=(2\pi)^{-n/2}\exp\left[-\frac12\sum_i(x_i-\theta)^2\right]
=h(x)\exp\left[-\frac n2(\bar x-\theta)^2\right],
\]
其中 $h(x)=(2\pi)^{-n/2}\exp[-\frac12\sum_i(x_i-\bar x)^2]$ 与 $\theta$ 无关。
\step{Step 2：先验假设概率与贝叶斯因子}
设 $P(H_i)=\pi_i$。两种假设均为单点，故
\begin{align*}
\BF_{01}(x)
&=\frac{h(x)e^{-n\bar x^2/2}}{h(x)e^{-n(\bar x-1)^2/2}}\\
&=\exp\left[-\frac n2\{\bar x^2-(\bar x-1)^2\}\right]
=\boxed{\exp\left[-\frac n2(2\bar x-1)\right].}
\end{align*}
所以 $\bar x<1/2$ 时数据相对支持 $H_0$，$\bar x>1/2$ 时相对支持 $H_1$，$\bar x=1/2$ 时贝叶斯因子为 $1$。

\step{Step 3：代入板书数值，并结合先验作决策}
当 $n=10,\bar x=2$ 时，
\[
\BF_{01}=e^{-15}\approx3.0590\times10^{-7},\qquad
\BF_{10}=e^{15}\approx3.2690\times10^6.
\]
若 $\pi_0=\pi_1$，则
\[
P(H_0\mid x)=\frac{e^{-15}}{1+e^{-15}}\approx3.0590\times10^{-7},
\qquad P(H_1\mid x)\approx0.9999996941.
\]
两种错误损失相同时选择 $H_1$。若希望选择 $H_0$，则须
\[
e^{-15}\frac{\pi_0}{\pi_1}>1
\quad\Longleftrightarrow\quad\frac{\pi_0}{\pi_1}>e^{15}.
\]
例如即使先验赔率达到 $10000:1$，后验赔率仍为
\[
\frac{\alpha_0}{\alpha_1}=10000e^{-15}\approx0.003059<1,
\]
仍选择 $H_1$。这说明足够强的样本证据可以改变原有的先验倾向。

\subsection{复合假设对复合假设：积分与最大化的区别}
当 $H_i$ 内含多个参数值时，贝叶斯因子采用\textbf{按条件先验加权积分}：
\[
\BF_{01}(x)=\frac{\int_{\Theta_0}p(x\mid\theta)g_0(\theta)\,d\theta}
{\int_{\Theta_1}p(x\mid\theta)g_1(\theta)\,d\theta}.
\]
板书还对照写出广义似然比的一种形式
\[
\lambda(x)=\frac{\sup_{\theta\in\Theta_0}p(x\mid\theta)}
{\sup_{\theta\in\Theta}p(x\mid\theta)}.
\]
这里分子、分母是最大值；贝叶斯因子的分子、分母是边缘似然。不能把积分直接替换为最大值。

\subsection{例2.4.3：正态后验下的复合假设检验}
设 $X_i\mid\theta\overset{iid}{\sim}N(\theta,1)$，先验 $\theta\sim N(0.5,2)$，$n=10,\bar x=1.5$，检验
\[
H_0:\theta\leq1,\qquad H_1:\theta>1.
\]
\step{Step 1：似然核}
\[
L(\theta;x)\propto\exp\left[-\frac n2(\theta-\bar x)^2\right].
\]
\step{Step 2：先验与先验假设概率}
\[
\pi(\theta)=\frac1{\sqrt{4\pi}}e^{-(\theta-0.5)^2/4},
\qquad \pi_0=\Phi\left(\frac{1-0.5}{\sqrt2}\right)\approx0.638163,
\]
所以 $\pi_1\approx0.361837$，先验赔率 $\pi_0/\pi_1\approx1.76366$。
\step{Step 3：后验核、配方及归一化}
\[
\pi(\theta\mid x)\propto\exp\left\{-\frac12
\left[n(\theta-\bar x)^2+\frac12(\theta-0.5)^2\right]\right\}.
\]
识别为 $N(\mu_n,S_n^2)$，其中
\[
S_n^2=\frac1{n+1/2}=\frac1{10.5},\qquad
\mu_n=\frac{n\bar x+0.5/2}{n+1/2}
=\frac{15.25}{10.5}\approx1.452381.
\]
\step{Step 4：后验假设概率}
\[
S_n\approx0.308607,\qquad
\alpha_0=\Phi\left(\frac{1-\mu_n}{S_n}\right)
=\Phi(-1.465882)\approx0.071340,
\]
\[
\alpha_1=1-\alpha_0\approx0.928660,
\qquad\frac{\alpha_0}{\alpha_1}\approx0.07682.
\]
\step{Step 5：贝叶斯因子及结论}
\[
\BF_{01}=\frac{\alpha_0/\alpha_1}{\pi_0/\pi_1}
\approx\frac{0.07682}{1.76366}=0.04356,\qquad
\BF_{10}\approx22.96.
\]
两种错误损失相同时选择 $H_1$。虽然先验更倾向 $H_0$，观察数据后的后验概率与贝叶斯因子都明显倾向 $H_1$。

\subsection{简单原假设对复合备择假设：点质量与连续先验的混合}
考虑 $H_0:\theta=\theta_0$ 与 $H_1:\theta\ne\theta_0$。若仅取连续先验，则单点的先验概率为零，无法赋予 $H_0$ 正的先验概率。板书的处理是另给原假设一个点质量：
\[
\pi(d\theta)=\pi_0\delta_{\theta_0}(d\theta)
+(1-\pi_0)g_1(\theta)\,d\theta,\qquad0<\pi_0<1,
\]
其中 $\delta_{\theta_0}$ 是在 $\theta_0$ 处的单位点质量，$g_1$ 是 $H_1$ 下的条件先验密度。

原假设与备择假设下的边缘似然分别为
\[
m_0(x)=p(x\mid\theta_0),\qquad
m_1(x)=\int_{\theta\ne\theta_0}p(x\mid\theta)g_1(\theta)\,d\theta,
\]
总边缘密度为 $m(x)=\pi_0m_0(x)+(1-\pi_0)m_1(x)$。所以
\[
\BF_{01}=\frac{p(x\mid\theta_0)}{m_1(x)},\qquad
\boxed{P(H_0\mid x)=\left[1+\frac{1-\pi_0}{\pi_0}\frac1{\BF_{01}}\right]^{-1}.}
\]
若 $\pi_0=1/2$，简化为
\[
P(H_0\mid x)=\frac{\BF_{01}}{1+\BF_{01}},\qquad
P(H_1\mid x)=\frac1{1+\BF_{01}}.
\]
连续密度 $g_1$ 在 $\theta_0$ 处的取值不影响积分；它与额外设置的点质量不能混淆。

\subsection{例2.4.4：二项概率等于二分之一的检验}
设 $X\mid\theta\sim\Bin(n,\theta)$，检验
\[
H_0:\theta=\tfrac12,\qquad H_1:\theta\ne\tfrac12.
\]
\step{Step 1：似然函数}
\[
p(x\mid\theta)=\binom nx\theta^x(1-\theta)^{n-x}.
\]
\step{Step 2：混合先验}
取 $P(H_0)=\pi_0$，$P(H_1)=1-\pi_0$，在 $H_1$ 下用 $(0,1)$ 上的均匀条件先验 $g_1(\theta)=1$。这里不能误写成整个先验只有 $U(0,1)$。
\step{Step 3：分别求边缘似然}
\[
m_0(x)=\binom nx2^{-n},
\]
\begin{align*}
m_1(x)&=\int_0^1\binom nx\theta^x(1-\theta)^{n-x}\,d\theta\\
&=\binom nx B(x+1,n-x+1)
=\frac{n!}{x!(n-x)!}\frac{x!(n-x)!}{(n+1)!}
=\frac1{n+1}.
\end{align*}
\step{Step 4：贝叶斯因子与后验概率}
\[
\boxed{\BF_{01}(x)=\frac{m_0(x)}{m_1(x)}
=\frac{(n+1)!}{2^n x!(n-x)!}.}
\]
当 $\pi_0=1/2,n=5,x=3$ 时，
\[
\BF_{01}=\frac{6!}{2^5 3!2!}=\frac{15}{8},\qquad
P(H_0\mid x)=\frac{15}{23},\quad P(H_1\mid x)=\frac8{23}.
\]
两种错误损失相同时选择 $H_0$。与例2.4.1相比，这次原假设是单点而非半个区间，先验构造与边缘似然也随之不同。

\section{预测}
\subsection{先验预测与后验预测}
设 $X$ 是已经观测的样本，$Z$ 是未来观测。给定 $\theta$ 后，假定 $Z$ 与 $X$ 条件独立，且 $Z$ 的条件密度或概率质量函数为 $g(z\mid\theta)$。

\textbf{尚未观察样本时}，用先验对未知参数平均，得到先验预测分布
\[
m(z)=\int g(z\mid\theta)\pi(\theta)\,d\theta.
\]
\textbf{已经观察到 $x$ 时}，把先验换成后验，得到后验预测分布
\[
\boxed{m(z\mid x)=\int g(z\mid\theta)\pi(\theta\mid x)\,d\theta
=E[g(z\mid\theta)\mid x].}
\]
老师强调：先求参数的后验，再对未来分布平均。一般不能把参数直接换成后验均值，因为
\[
E[g(z\mid\theta)\mid x]\ne g\{z\mid E(\theta\mid x)\}
\]
通常成立。代入一个参数估计值，会漏掉参数本身的后验不确定性。

\subsection{预测均值、方差与点预测}
由条件期望公式，
\[
\boxed{E(Z\mid x)=E\{E(Z\mid\theta)\mid x\}.}
\]
由条件方差分解，
\[
\boxed{\Var(Z\mid x)=E\{\Var(Z\mid\theta)\mid x\}
+\Var\{E(Z\mid\theta)\mid x\}.}
\]
第一项来自给定参数后未来观测自身的随机波动，第二项来自参数的后验不确定性。平方损失下，最佳点预测为 $\hat z=E(Z\mid x)$；绝对损失下，最佳点预测为后验预测分布的中位数。

\subsection{例2.5.1：二项模型的 Beta--二项预测分布}
过去 $n$ 次试验中的成功次数为 $X$，未来 $k$ 次试验的成功次数为 $Z$。在给定 $\theta$ 时，
\[
X\mid\theta\sim\Bin(n,\theta),\qquad
Z\mid\theta\sim\Bin(k,\theta),\qquad X\perp Z\mid\theta.
\]
\step{Step 1：已有样本的似然函数}
\[
L(\theta;x)=\binom nx\theta^x(1-\theta)^{n-x}.
\]
\step{Step 2：先验密度}
\[
\pi(\theta)=\frac1{B(\alpha,\beta)}
\theta^{\alpha-1}(1-\theta)^{\beta-1},\quad0<\theta<1.
\]
\step{Step 3：参数的后验分布}
相乘后识别为
\[
\theta\mid x\sim\Be(a,b),\qquad a=\alpha+x,\quad b=\beta+n-x.
\]
\step{Step 4：对后验积分，得到未来次数的预测分布}
对 $z=0,1,\ldots,k$，
\begin{align*}
P(Z=z\mid x)
&=\int_0^1\binom kz\theta^z(1-\theta)^{k-z}
\frac{\theta^{a-1}(1-\theta)^{b-1}}{B(a,b)}\,d\theta\\
&=\frac{\binom kz}{B(a,b)}
\int_0^1\theta^{a+z-1}(1-\theta)^{b+k-z-1}\,d\theta\\
&=\boxed{\binom kz\frac{B(a+z,b+k-z)}{B(a,b)}}.
\end{align*}
这称为 Beta--二项预测分布。用 Gamma 函数展开为
\[
P(Z=z\mid x)=\binom kz
\frac{\Gamma(a+b)\Gamma(a+z)\Gamma(b+k-z)}
{\Gamma(a)\Gamma(b)\Gamma(a+b+k)}.
\]
\step{Step 5：代入板书数据 $n=10,x=3,k=5$}
取均匀先验 $\alpha=\beta=1$，则 $a=4,b=8$，即
\[
\theta\mid x=3\sim\Be(4,8).
\]
例如
\[
P(Z=0\mid x=3)=\frac{B(4,13)}{B(4,8)}
=\frac{\Gamma(12)\Gamma(13)}{\Gamma(8)\Gamma(17)}
=\frac{33}{182}\approx0.1813.
\]
其余各项同理得到
\begin{center}
\begin{tabular}{c*{6}{c}}
\toprule
$z$&0&1&2&3&4&5\\
\midrule
$P(Z=z\mid x=3)$&$99/546$&$165/546$&$150/546$&$90/546$&$35/546$&$7/546$\\
\bottomrule
\end{tabular}
\end{center}
这六个概率之和为 $1$。从表中可见，最可能的未来成功次数是 $1$；平方损失下的点预测却是
\[
\hat z=E(Z\mid x)=E(k\theta\mid x)
=k\frac a{a+b}=5\times\frac4{12}=\frac53.
\]
点预测为 $5/3$ 不要求未来次数实际能取 $5/3$；它是平方损失下的最优实数行动值。

板书还计算了预测区间 $[0,3]$ 的概率：
\[
P(0\leq Z\leq3\mid x=3)
=\frac{99+165+150+90}{546}=\frac{12}{13}\approx0.9231.
\]
因此它可作为可信水平至少 $90\%$ 的预测区间。

\step{补充：若没有过去样本，均匀先验的预测是什么？}
直接对先验积分：
\[
P(Z=z)=\binom kz\int_0^1\theta^z(1-\theta)^{k-z}\,d\theta
=\binom kz B(z+1,k-z+1)=\frac1{k+1}.
\]
所以 $Z$ 在 $0,1,\ldots,k$ 上均匀分布。直接代入均匀先验的均值 $1/2$，会得到 $\Bin(k,1/2)$；当 $k\geq2$ 时，它与这里的先验预测分布不同。

\subsection{正态预测及例2.5.2的两架天平}
设第一架天平测量 $n$ 次，$X_i\mid\theta\overset{iid}{\sim}N(\theta,\sigma^2)$，先验为 $N(\mu,\tau^2)$。按前面的似然乘先验并配方后，
\[
\theta\mid x\sim N(\mu_n,\tau_n^2),\qquad
\tau_n^2=\left(\frac n{\sigma^2}+\frac1{\tau^2}\right)^{-1},\quad
\mu_n=\tau_n^2\left(\frac{n\bar x}{\sigma^2}+\frac\mu{\tau^2}\right).
\]
第二架天平的未来称量值为 $Z\mid\theta\sim N(\theta,\sigma_2^2)$。为简化配方，令 $v=\tau_n^2$，$w=\sigma_2^2$，则
\begin{align*}
m(z\mid x)
&=\frac1{2\pi\sqrt{vw}}\int_{-\infty}^{\infty}
\exp\left\{-\frac12\left[\frac{(z-\theta)^2}{w}
+\frac{(\theta-\mu_n)^2}{v}\right]\right\}\,d\theta.
\end{align*}
沿用课件的配方过程，记
\[
A=\frac1w+\frac1v,\qquad B=\frac zw+\frac{\mu_n}v,\qquad
C=\frac{z^2}w+\frac{\mu_n^2}v.
\]
则指数中的二次式为
\[
A\theta^2-2B\theta+C
=A\left(\theta-\frac BA\right)^2+C-\frac{B^2}A,
\]
其中 $C-B^2/A=(z-\mu_n)^2/(v+w)$。利用正态核的积分为 $\sqrt{2\pi/A}$，得到
\[
m(z\mid x)=\frac1{\sqrt{2\pi(v+w)}}
\exp\left[-\frac{(z-\mu_n)^2}{2(v+w)}\right].
\]
故
\[
\boxed{Z\mid x\sim N(\mu_n,\sigma_2^2+\tau_n^2).}
\]
也可用板书中的条件矩公式核对：
\[
E(Z\mid x)=E(\theta\mid x)=\mu_n,
\]
\[
\Var(Z\mid x)=E(\sigma_2^2\mid x)+\Var(\theta\mid x)
=\sigma_2^2+\tau_n^2.
\]

\step{参数可信区间与未来观测预测区间的区别}
在同一可信水平 $95\%$ 下，
\[
\begin{array}{ll}
\text{参数 }\theta\text{ 的可信区间：}&\mu_n\pm1.96\tau_n,\\[3pt]
\text{未来观测 }Z\text{ 的预测区间：}&\mu_n\pm1.96\sqrt{\sigma_2^2+\tau_n^2}.
\end{array}
\]
后者还包括未来观测本身的误差，因此更宽。板书中若统一把未来观测方差记为 $\sigma^2$，只需在上式以 $\sigma^2$ 替换 $\sigma_2^2$。

\section{似然原理}
\subsection{似然函数的含义与比例关系}
给定模型 $p(x\mid\theta)$，同一表达式可以从两个角度看：固定 $\theta$、让 $x$ 变化时，是样本的抽样分布；固定已经观测到的 $x$、让 $\theta$ 变化时，是似然函数
\[
L(\theta;x):=p(x\mid\theta).
\]
似然\textbf{不是参数的后验密度}，不要求对 $\theta$ 的积分等于 $1$。先验乘似然后，才通过贝叶斯公式得到后验。

若两次试验分别得到观测 $x,y$，并且对所有 $\theta$ 有
\[
L_1(\theta;x)=c\,L_2(\theta;y),\qquad c>0\text{ 且与 }\theta\text{ 无关},
\]
则两者对任意两个参数值的相对支持相同：
\[
\frac{L_1(\theta_a;x)}{L_1(\theta_b;x)}
=\frac{L_2(\theta_a;y)}{L_2(\theta_b;y)}.
\]
\textbf{似然原理}主张：对参数的统计推断应依据实际观测给出的似然；比例相同的似然提供相同的相对证据。

取相同先验且后验可归一化时，比例常数在贝叶斯公式中抵消：
\[
\pi_1(\theta\mid x)
=\frac{cL_2(\theta;y)\pi(\theta)}{c\int L_2(t;y)\pi(t)\,dt}
=\pi_2(\theta\mid y).
\]
因此此时得到相同的后验及基于该后验的推断。

\subsection{课堂引例：固定次数与达到成功数后停止}
\textbf{方案A：}预先规定做 $10$ 次试验，观察到 $3$ 次成功。
\[
L_A(\theta)=\binom{10}{3}\theta^3(1-\theta)^7.
\]
\textbf{方案B：}一直做到第 $3$ 次成功，恰在第 $10$ 次试验停止。最后一次必须成功，前 $9$ 次有 $2$ 次成功，因此
\[
L_B(\theta)=\binom92\theta^3(1-\theta)^7.
\]
两者之比为
\[
\frac{L_A(\theta)}{L_B(\theta)}=\frac{\binom{10}{3}}{\binom92}
=\frac{10}{3},
\]
与 $\theta$ 无关。若均取 $\Be(\alpha,\beta)$ 先验，则两个方案都得到
\[
\theta\mid\text{data}\sim\Be(\alpha+3,\beta+7).
\]
停止规则不同，但这两组实际观测给出的似然成比例，采用同一先验时后验相同。

\subsection{例2.6.1：抛硬币的两个方案}
设正面概率为 $\theta$，考虑
\[
H_0:\theta=\tfrac12,\qquad H_1:\theta>\tfrac12.
\]
\step{第一步：两种试验的似然函数}
方案一预先规定抛 $12$ 次，得到 $9$ 次正面、$3$ 次反面：
\[
L_1(\theta)=\binom{12}{9}\theta^9(1-\theta)^3
=220\theta^9(1-\theta)^3.
\]
方案二一直抛到出现第 $3$ 次反面停止，此前共出现 $9$ 次正面。最后一次是反面，前 $11$ 次恰有 $9$ 次正面、$2$ 次反面：
\[
L_2(\theta)=\binom{11}{9}\theta^9(1-\theta)^3
=55\theta^9(1-\theta)^3.
\]
故 $L_1=4L_2$，满足似然成比例的条件。

\step{第二步：频率方法的两个 $p$ 值}
在方案一中，$H_0$ 下 $X\sim\Bin(12,1/2)$，所以
\[
p_1=P_{H_0}(X\geq9)
=\frac{\binom{12}{9}+\binom{12}{10}+\binom{12}{11}+\binom{12}{12}}{2^{12}}
=\frac{299}{4096}\approx0.07300.
\]
方案二令 $Y$ 表示第 $3$ 次反面出现之前的正面次数，则
\[
P_\theta(Y=j)=\binom{j+2}{j}\theta^j(1-\theta)^3,\qquad j=0,1,\ldots,
\]
因此
\[
p_2=P_{H_0}(Y\geq9)
=\sum_{j=9}^{\infty}\binom{j+2}{j}2^{-(j+3)}.
\]
“第3次反面之前至少有9次正面”等价于“前11次中至少有9次正面”，故也可以算作
\[
p_2=\frac{\binom{11}{9}+\binom{11}{10}+\binom{11}{11}}{2^{11}}
=\frac{67}{2048}\approx0.03271.
\]
在 $5\%$ 显著性水平下，方案一不拒绝 $H_0$，方案二拒绝 $H_0$。这是因为两个 $p$ 值依据不同抽样方案下的尾部事件计算。

\step{第三步：先验假设概率及备择下的条件先验}
按板书取
\[
P(H_0)=P(H_1)=\frac12,\qquad
\theta\mid H_1\sim U(1/2,1),
\]
所以 $g_1(\theta)=2\ind_{(1/2,1)}(\theta)$。

\step{第四步：边缘似然与贝叶斯因子}
把方案 $i$ 的似然统一写为 $L_i(\theta)=k_i\theta^9(1-\theta)^3$，其中 $k_1=220,k_2=55$。则
\[
m_{0,i}=k_i2^{-12},\qquad
m_{1,i}=2k_i\int_{1/2}^1\theta^9(1-\theta)^3\,d\theta.
\]
相除后 $k_i$ 抵消，故两个方案的贝叶斯因子相同：
\[
\BF_{01}=\frac{2^{-12}}{2\displaystyle\int_{1/2}^1\theta^9(1-\theta)^3\,d\theta}.
\]
沿用板书的多项式展开，
\[
\theta^9(1-\theta)^3=\theta^9-3\theta^{10}+3\theta^{11}-\theta^{12},
\]
\[
I=\int_{1/2}^1\theta^9(1-\theta)^3\,d\theta
=\left[\frac{\theta^{10}}{10}-\frac{3\theta^{11}}{11}
+\frac{\theta^{12}}4-\frac{\theta^{13}}{13}\right]_{1/2}^1
\approx0.0003335166.
\]
从而
\[
\BF_{01}=\frac{1430}{3907}\approx0.36601,\qquad
\BF_{10}\approx2.73217.
\]

\step{第五步：后验概率与结论}
先验赔率为 $1$，所以两个方案都给出
\[
P(H_0\mid\text{data})=\frac{\BF_{01}}{1+\BF_{01}}
=\frac{1430}{5337}\approx0.26794,
\]
\[
P(H_1\mid\text{data})=\frac1{1+\BF_{01}}
=\frac{3907}{5337}\approx0.73206.
\]
两种错误损失相同时，两方案均选择 $H_1$。例子说明：似然成比例并取相同先验时，贝叶斯计算不因这两种停止规则的差别而改变；频率 $p$ 值则可能不同。

\end{document}
